\honest{Words as Mental Paintbrush Handles}{Eliezer Yudkowsky}{2008}

I am writing from the AGI-08 conference, and I think we are done with the mathematical posts, at least for now.

Suppose I tell you that the lamps in this hotel have triangular lightbulbs. Picture one. When I did, ``as best as my introspection could determine,'' I first saw a pyramid with sharp edges, then the edges were smoothed, and then I saw a loop of fluorescent tube bent into a triangle. As far as I can tell, no words were involved.

For some decades, I say, researchers debated whether people really have mental images or only think they do, while ``actually just having a little `chair' label, like a LISP token.'' I try hard not to call this ``spectacularly silly,'' because of the hindsight effect, and then I call it that. \nb{A few lines later the post grants that the question ``used to be harder to answer, back in the day of the controversy.''}

The doubting side, I think, was ``mostly a deranged legacy of behaviorism, which denied the existence of thoughts in humans.'' \nb{The doubting side, whose main advocate was Zenon Pylyshyn, held that images are stored as symbolic internal representations, and traditional behaviorism tried to keep internal representations out of psychology. Skinner, according to the \textsc{Stanford Encyclopedia of Philosophy}, ``readily admits that private thoughts and so on exist.''}

How do we know images are real? In the mental-rotation experiments, the time needed to match two objects grows with the angle between them. ``Today we can actually neuroimage the little pictures in the visual cortex,'' so the brain ``really does represent a detailed image.'' I cite Stephen Kosslyn's book, subtitled ``The Resolution of the Imagery Debate.'' \nb{The picture view is the one the \textsc{Stanford Encyclopedia} now favors. In 2003, though, Pylyshyn was still arguing that the neuroimaging evidence ``does not address the format of mental images,'' and in 2015 Kosslyn co-wrote a paper subtitled ``Ending the imagery debate.''}

Now my thesis. People get into trouble with words partly because they do not see how much lies behind them. ``Apple'' is five letters, but the concept behind it can paint a picture in your visual cortex, combine with ``cheese,'' recognize an apple, taste one in a pie, and perhaps guide your hand to eat one. Words are paintbrushes: ``you can use them to draw images in your own mind,'' and in other people's. \nb{Some people report no mental imagery at all, and they use words normally. The other abilities in the list do not depend on pictures.}

So do not make the behaviorists' mistake. A word is a pointer, and sooner or later you must look where it points. What does a word point to?
